%\section*{Partie D - Synthèse d'une loi de commande de l'actionneur}
\section{Synthèse d'une loi de commande de l'actionneur}

\begin{obj}
Proposer une loi de commande de l'actionneur mettant en mouvement la nacelle et vérifier les exigences relatives à la sécurité, à l'accessibilité, au panorama et au confort des passagers.
\end{obj}

La partie précédente a montré l'instabilité d'une commande en couple de l'actionneur. Pour contrôler la vitesse de la nacelle $N$, il est nécessaire de stabiliser le comportement dynamique de l'ensemble $\Sigma$ régi par l'équation Eq. 2. L'actionneur, délivrant le couple $c_{e m}(t)$, est commandé en couple par l'intermédiaire du préactionneur par la grandeur $c_{\text {ref }}(t)$.

\subsection*{Rappels et notations}
\begin{itemize}
  \item $\omega_{m c}(t)$ représente la vitesse angulaire de consigne du rotor de l'actionneur, exprimée en $\operatorname{rad} \cdot \mathrm{s}^{-1} . \Omega_{m c}(p)$ est sa transformée de Laplace ;
  \item $\omega_{m}(t)$ représente la vitesse angulaire du rotor de l'actionneur, exprimée en $\mathrm{rad} \cdot \mathrm{s}^{-1} . \Omega_{m}(p)$ est sa transformée de Laplace;
  \item $\theta_{m}(t)$ représente la position angulaire du rotor de l'actionneur, exprimée en $\operatorname{rad} . \Theta_{m}(p)$ est sa transformée de Laplace;
  \item $c_{e m}(t)$ est l'action mécanique du type couple générée par l'actionneur sur le réducteur de vitesse exprimée en $\mathrm{N} \cdot \mathrm{m}$;
  \item $c_{\text {ref }}(t)$ est la grandeur de commande de l'ensemble \{Préactionneur + actionneur\} exprimée en $\mathrm{N} \cdot \mathrm{m}$. Le temps de réponse de cet ensemble étant faible comparativement à la dynamique du système étudié, on considère $c_{e m}(t)=c_{r e f}(t)$ et $\Delta c_{e m}(t)=\Delta c_{\text {ref }}(t)$;
  \item $c_{R}(t)$ représente une action mécanique résistante rapportée sur l'axe de rotation du rotor de l'actionneur, exprimée en $\mathrm{N} \cdot \mathrm{m} . C_{R}(p)$ est sa transformée de Laplace.
\end{itemize}

%\section*{I - Stabilisation du comportement dynamique de l'ensemble de mise en mouvement de la nacelle}
\subsection{Stabilisation du comportement dynamique de l'ensemble de mise en mouvement de la nacelle}

\begin{obj}
Élaborer une commande de l'actionneur assurant la stabilité de la dynamique de l'ensemble de mise en mouvement de la nacelle.
\end{obj}

Pour contrôler la vitesse linéaire $v(t)$ de la nacelle $N$, il est nécessaire de stabiliser le comportement dynamique de l'ensemble $\Sigma$ régi par l'équation Eq.2. À cet effet, la variation de couple $\Delta c_{\text {ref }}(t)=\Delta c_{e m}(t)$ est commandée par un retour informationnel fonction de la vitesse angulaire de consigne et de la vitesse angulaire du rotor de l'actionneur telle que :

$$
\Delta c_{e m}(t)=-K_{1} \cdot \omega_{m}(t)-K_{2} \cdot \int_{0}^{t}\left(\omega_{m}(x)-K_{3} \cdot \omega_{m c}(x)\right) \cdot \mathrm{d} x \quad \text { Eq. } 3
$$

où $K_{1}, K_{2}$ et $K_{3}$ sont des constantes positives à déterminer.\\

%Q22.
\question{À l'aide des équations Eq. 2, Eq. 3, en l'absence de perturbation $\left(C_{R}(p)=0\right)$ et en supposant les conditions de Heaviside vérifiées, déterminer l'expression de la fonction de transfert $H_{\text {stab }}(p)$ en fonction de $K_{1}, K_{2}, K_{3}$, $J_{0}, k_{c r z}$ et $k_{c i n}$ telle que :
$\Omega_{m}(p)=\indice{H}{stab}(p) \cdot \Omega_{m c}(p)$.}


Afin de limiter les oscillations des câbles et leurs sollicitations en traction notamment, le comportement souhaité de la boucle de stabilisation est telle que $H_{\text{stab}}(p)$ s'écrive sous la forme :

$$
H_{sta b}(p)=\frac{1}{1+\frac{2 \xi_{s}}{\omega_{0 s}} \cdot p^{+} \frac{p^{2}}{\omega_{0 s}^{2}}} .
$$

Le cahier des charges relatif à la boucle de stabilisation est fourni dans le Tableau 4.

\begin{table}[h]
\begin{center}
\begin{tabular}[t]{|l|l|l|}
\hline
Performances & Critères & Niveaux \\
\hline
Stabilité & Facteur d'amortissement de la boucle de stabilisation & $\xi_{s} \geqslant 1$ \\
\hline
Rapidité & Pulsation propre du système non amorti de la boucle de stabilisation & $\omega_{0 s}=3 \mathrm{rad} \cdot \mathrm{s}^{-1}$ \\
\hline
\end{tabular}
\captionsetup{labelformat=empty}
\caption{Tableau 4 - Cahier des charges de la boucle de stabilisation}
\end{center}
\end{table}

%Q23. 
\question{Déterminer les expressions littérales de $K_{1}, K_{2}$ et $K_{3}$ en fonction de $\xi_{s}, \omega_{0 s}, J_{0}, k_{c r z}$ et $k_{c i n}$.}

Indépendamment des résultats précédents, les valeurs des paramètres du modèle (Eq. 2) adoptées sont :

$$
J_{0}=16,8 \mathrm{~kg} \cdot \mathrm{~m}^{2} \quad ; \quad k_{\text {cin }}=\frac{1}{204} \mathrm{~m} \cdot \mathrm{rad}^{-1} \quad ; \quad k_{\text {cr } z}=4,5 \mathrm{~N} .
$$

%Q24. 
\question{En déduire les conditions sur les valeurs numériques des paramètres $K_{1}$ et $K_{2}$ permettant de satisfaire les critères du Tableau 4.}

En prenant en compte la perturbation $C_{R}(p)$, il est possible de démontrer que la fonction de transfert $G_{s}(p)=\frac{\Omega_{m}(p)}{C_{R}(p)}$ s'écrit sous la forme $G_{s}(p)=\frac{\Omega_{m}(p)}{C_{R}(p)}=-H_{\text {stab }}(p) \cdot \frac{p}{J_{0} \omega_{o s}^{2}}$.\\


Soit $n$ le nombre de passagers embarqués dans la nacelle. Considérons $c_{R}(t)=C_{R 0}(n)=$ constante.\\

%Q25. 
\question{Montrer que la masse des $n$ passagers dans la nacelle a une influence sur la position angulaire du moteur en régime permanent et en conséquence sur la position de la nacelle $z(t)$.}


%\section*{II - Étude de la boucle de maîtrise des grandeurs cinématiques de l'actionneur}
\subsection{Étude de la boucle de maîtrise des grandeurs cinématiques de l'actionneur}

%Objectif\\
\begin{obj}
Déterminer les paramètres d'un correcteur permettant de satisfaire les exigences de sécurité et de confort des passagers.
\end{obj}

La dynamique du mécanisme de mise en mouvement de la nacelle étant désormais stabilisée, la vitesse $v(t)$ de la nacelle doit suivre un profil établi en fonction des critères de sécurité et de confort des passagers, de la durée de la visite et des positions de la nacelle à l'arrivée aux niveaux inférieur et supérieur. À cet effet, une boucle d'asservissement de la vitesse angulaire de l'actionneur est réalisée, conformément à la Figure 7. $C(p)$ représente le correcteur de la boucle d'asservissement de vitesse angulaire devant permettre de maîtriser les grandeurs cinématiques de la nacelle.

\begin{figure}[h]
\begin{center}
  \includegraphics[alt={},max width=\textwidth]{cd8229cf-9637-40cc-9aaf-03c75c002f2c-12_443_849_1025_616}
\captionsetup{labelformat=empty}
\caption{Figure 7 - Schéma-blocs simplifié de l'asservissement de vitesse angulaire de l'actionneur}
\end{center}
\end{figure}

Pour limiter les oscillations du câble, ainsi que les variations de la vitesse dues aux changements de pentes de la consigne de vitesse angulaire $\omega_{m c}(t)$, l'asservissement de vitesse angulaire du rotor de l'actionneur doit respecter les performances du Tableau 5.

\begin{table}[h]
\begin{center}
\begin{tabular}[t]{|l|l|l|}
\hline
Performances & Critères & Niveaux \\
\hline
\multirow[t]{3}{*}{Stabilité} & Dépassement relatif du $k^{\mathrm{e}}$ dépassement pour une entrée de consigne en échelon, noté $D_{k}^{\%}$ & $D_{k}^{\%}=0 \forall k$ \\
\hline
 & Marge de phase, notée $M \varphi$ & $M \varphi=80^{\circ}$ \\
\hline
 & Marge de gain, notée $M G$ & $M G \geqslant 50 \mathrm{~dB}$ \\
\hline
\multirow[t]{2}{*}{Précision} & Erreur en régime permanent sur la vitesse angulaire de l'actionneur, notée $\mu_{v}$, pour une entrée de consigne en échelon & $\mu_{v}=0$ \\
\hline
 & Sensibilité en régime permanent à une perturbation en échelon & Insensible \\
\hline
\end{tabular}
\captionsetup{labelformat=empty}
\caption{Tableau 5 - Cahier des charges de la boucle d'asservissement de vitesse angulaire}
\end{center}
\end{table}

La stabilisation étudiée dans la Partie D-I, et par le choix d'un facteur d'amortissement $\xi_{s} \geqslant 1$, permettent de considérer que la fonction de transfert de la boucle de stabilisation s'écrit :

$$
H_{s t a b}(p)=\frac{1}{(1+1,24 \cdot p) \cdot(1+0,09 \cdot p)} .
$$

Les ingénieurs proposent un choix de correcteur $C(p)$ du type Proportionnel Intégral, de fonction de transfert :

$$
C(p)=K_{i} \cdot \frac{1+T_{i} \cdot p}{T_{i} \cdot p} .
$$

La détermination du paramètre $T_{i}$ est telle qu'il réalise la compensation du pôle dominant de $H_{\text {stab }}(p)$.

\begin{itemize}
  \item[Q26.] Déterminer les deux pôles de $H_{\text {stab }}(p)$. Quel est le pôle dominant et justifier la réponse? En déduire la valeur numérique de $T_{i}$.
\end{itemize}

La Figure 8 fournit le diagramme de Bode de la fonction de transfert en boucle ouverte corrigée pour $K_{i}=1$ et la compensation du pôle dominant.

\begin{figure}[h]
\begin{center}
  \includegraphics[alt={},max width=\textwidth]{cd8229cf-9637-40cc-9aaf-03c75c002f2c-13_907_1389_377_338}
\captionsetup{labelformat=empty}
\caption{Figure 8 - Diagramme de Bode de la fonction de transfert en boucle ouverte corrigée pour $K_{i}=1$ et compensation du pôle dominant}
\end{center}
\end{figure}

%\begin{itemize}
 % \item[Q27.] 
%\end{itemize}


 % Q27 
\question{Déterminer numériquement la condition sur le paramètre $K_{i}$ permettant de satisfaire simultanément les critères de stabilité relatifs aux marges de gain et de phase.}


Par calculs ou par manipulations du schéma-blocs de la Figure 7, il est possible de démontrer que la fonction de transfert $G(p)=\frac{\Omega_{m}(p)}{C_{R}(p)}$ s'écrit sous la forme :

$$
G(p)=\frac{\Omega_{m}(p)}{C_{R}(p)}=\frac{a_{0} \cdot p^{2}}{b_{0}+b_{1} \cdot p+b_{2} \cdot p^{2}+b_{3} \cdot p^{3}} .
$$

Soit $n$ le nombre de passagers embarqués dans la nacelle. Considérons $c_{R}(t)=C_{R 0}(n)=$ constante.

\begin{itemize}
  \item[Q28.] Montrer que, quelque soit $n$, la masse de $n$ passagers dans la nacelle n'a pas d'influence sur la vitesse et la position de la nacelle en régime permanent.
\end{itemize}

La Figure 9 fournit le diagramme de Bode de la fonction de transfert en boucle fermée $F T B F(p)=\frac{\Omega_{m}(p)}{\Omega_{m c}(p)}$ et pour les valeurs pertinentes des paramètres du correcteur $C(p)$.


%Q29.
\question{Conclure quant à la satisfaction ou non de l'intégralité des performances attendues du cahier des charges du Tableau 5.}

Les études précédentes ont permis de déterminer une modélisation de la commande de l'actionneur afin de contrôler les grandeurs cinématiques. Il est donc désormais possible de vérifier les exigences liées à la sécurité, au panorama et au confort des passagers pour le cycle de vitesse souhaité décrit en partie A-I.

\begin{figure}[h]
\begin{center}
  \includegraphics[alt={},max width=\textwidth]{cd8229cf-9637-40cc-9aaf-03c75c002f2c-14_921_1396_121_336}
\captionsetup{labelformat=empty}
\caption{Figure 9 - Diagramme de Bode de $F T B F(p)$}
\end{center}
\end{figure}

%\section*{III - Vérification des performances liées à la sécurité et au confort des passagers}
\subsection{Vérification des performances liées à la sécurité et au confort des passagers}
%Objectif\\

\begin{obj}
Caractériser les écarts entre les performances attendues liées à la sécurité et au confort des passagers et les performances simulées (hors critères d'accessibilité de la nacelle et de panorama définis Partie A-I).
\end{obj}

L'ensemble des paramètres du modèle de la Figure 4 et les grandeurs d'entrées sont à présents connus. On considère le cycle ascensionnel de vitesse souhaitée de la nacelle $N$ établi en partie A-I, en absence de passager $n=0$ et en présence de $n=200$ passagers dans la nacelle.

Un frein mécanique permet de maintenir à l'arrêt la nacelle (en $z=0 \mathrm{~m}$ et à l'altitude maximale). On considère qu'à l'instant $t=0$ de la simulation, l'actionneur délivre un couple $c_{\text {em }}(t)$ tel que $c_{\text {em }}(t=0)=C_{R 0}(n)$ où $n$ représente le nombre de passagers embarqués dans la nacelle.\\
La simulation du modèle établi dans ce sujet permet d'obtenir les Figures 10 et 11 qui décrivent les évolutions des grandeurs cinématiques de la nacelle $N$.

\begin{figure}[h]
\begin{center}
  \includegraphics[alt={},max width=\textwidth]{cd8229cf-9637-40cc-9aaf-03c75c002f2c-14_795_1240_1864_413}
\captionsetup{labelformat=empty}
\caption{Figure 10 - Évolution des accélérations linéaires de la nacelle en fonction du nombre de passagers $n$}
\end{center}
\end{figure}

\begin{figure}[h]
\begin{center}
  \includegraphics[alt={},max width=\textwidth]{cd8229cf-9637-40cc-9aaf-03c75c002f2c-15_816_1240_121_413}
\captionsetup{labelformat=empty}
\caption{Figure 11 - Évolution des vitesses linéaires de la nacelle en fonction du nombre de passagers $n$}
\end{center}
\end{figure}

%Q30. 
\question{Conclure quant à la satisfaction ou non des critères des exigences liées à la sécurité et au confort des passagers (Tableau 2), quelque soit le nombre de passagers embarqués dans la nacelle $N$.}

